\documentclass[10pt,a4paper]{amsart}
\usepackage[margin=19mm]{geometry}
\usepackage{amsmath,amssymb,mathtools}
\usepackage[scr=rsfs,cal=euler]{mathalfa}
\usepackage{xfrac}
\usepackage[unicode,hidelinks,bookmarks=false]{hyperref}
\newcommand{\ie}{i.e.\ }
\newcommand{\cf}{cf.\ }
\newcommand{\resp}{respectively\ }
\newcommand{\Subsection}{\subsection}
\providecommand{\nobreakdash}{\nobreak\hbox{-}}
\setlength{\emergencystretch}{2em}
\multlinegap0pt
\usepackage{enumerate}
\usepackage{csquotes}
\usepackage[shortlabels]{enumitem}
\usepackage{amssymb}
\usepackage{mathtools}
\newtheorem{assumption}{Assumption}
\newtheorem{proposition}{Proposition}
\newcommand{\N}{\mathbb{N}}
\newcommand{\R}{\mathbb{R}}
\newcommand{\YY}{X}
\title{Evolution variational inequalities with general costs}
\author{Pierre-Cyril Aubin-Frankowski, Giacomo Enrico Sodini, Ulisse Stefanelli}
\date{Source: arXiv:2505.00559v2 (CC BY 4.0)}
\begin{document}
\maketitle

\noindent\emph{Source and licence.} Evolution variational inequalities with general costs. Version arXiv:2505.00559v2 (CC BY 4.0), distributed by arXiv under the Creative Commons Attribution 4.0 International licence, http://creativecommons.org/licenses/by/4.0/, as recorded in the arXiv licence metadata for this exact version and shown on the abstract page https://arxiv.org/abs/2505.00559v2. Full licence notice: \texttt{licenses/arxiv-2505.00559-CC-BY-4.0.txt}.\par\smallskip

\noindent\textit{Editorial scope.} Counted unit: the article's proposition prop:cauchy (source0.tex lines 1389--1397). The article's complete original proof of it is delivered with it (source0.tex lines 1399--1428, one \texttt{proof} environment of 30 lines). No mathematics is written, repaired or re-derived: the statement and the proof are the article's own lines, in the article's own notation, and no retained line is substituted or re-flowed.  Retained uncounted as the source-local context the proof needs: lines 276--278; lines 1304--1309; lines 1356--1361.  Nothing was omitted from any retained span, so the package carries no omission record.  No retained line contains a citation, so the wrapper carries no bibliography and no bibliography entry.  No retained line contains a figure environment, an image-inclusion command or drawing code, so no image is delivered and the package declares no asset dependency.  Editorial formatting only, in addition: an AMS \texttt{amsart} preamble that restates the article's own declarations and macros used by the retained lines, namely the environments \texttt{assumption}, \texttt{proposition} on separate counters, together with the standard packages those lines need.  The retained environments print in this document as Assumption 1, Proposition 1 in order of appearance, whereas the article prints the same items with the numbering of its own full document.  The complete native source of the article as distributed in the arXiv e-print for this exact version is bundled unchanged as \texttt{sources/177423/source0.tex}.
    \begin{equation}\label{eq:addass}\tag{Diss}
     c \ge 0 \quad \text{ and } \quad c(x,x')=0 \text{ if and only if } x=x'.
    \end{equation}

\begin{assumption}\label{ass:ass1}  Let  $c:X \times \YY\to \R$ and $g: X \to \R \cup {+\infty}$  be given  and  assume   that there exist $\bar{\tau}>0$ and $\lambda \ge 0$ such that
\begin{enumerate}
\item $-g$ is $\lambda \tau$-strongly $c/\tau$-cross-concave for every $0<\tau< \bar{\tau}$;
\item for every $x_0 \in X$ and every $0< \tau<\bar{\tau}$ there exists some $y_0 \in \mathcal{R}_c(x_0)$ such that $\mathcal{P}_{c/\tau}(-g,y_0) \ne \emptyset$.
\end{enumerate}
\end{assumption}

    \begin{align}
&[c(\bar{x}_{t}^{\tau/2^p},\bar{y}_{t}^\tau)+c(\bar{x}_{t}^\tau,\bar{y}_{t}^{\tau/2^p})]\nonumber\\
 &  \qquad +\sum_{i=0}^{n-1}[c(\bar{x}_{i\tau}^{\tau/2^p},\bar{y}_{(i+1)\tau}^\tau)- c(\bar{x}_{(i+1)\tau}^{\tau/2^p},\bar{y}_{i\tau}^\tau) + c(\bar{x}_{i\tau}^\tau,\bar{y}_{(i+1)\tau}^{\tau/2^p})-c(\bar{x}_{(i+1)\tau}^\tau,\bar{y}_{i\tau}^{\tau/2^p})]\nonumber\\ 
 & \qquad  +2\sum_{i=0}^{n-1}c(\bar{x}_{(i+1)\tau}^\tau, \bar{y}_{i\tau}^\tau)+ 2\sum_{j=0}^{n2^p-1}c(\bar{x}_{(j+1)\tau/2^p}^{\tau/2^p}, \bar{y}_{j\tau/2^p}^{\tau/2^p})\nonumber\\
    & \quad \le \tau(2g(x_0)-g(\bar{x}_{t}^\tau)-g(\bar{x}_{t}^{\tau/2^p})).\label{eq:Cauchy_ineq_full}
\end{align}

\begin{proposition}[Cauchy estimate]\label{prop:cauchy}  Let $c$ be a cost satisfying \eqref{eq:addass}  and let  Assumption \ref{ass:ass1} hold. Assume additionally  that $c$     decomposes as $c=c_1+c_2$ with 
\begin{enumerate}
    \item  $c_1\geq 0$  symmetric: $c_1(x_1,x_2)=c_1(x_2,x_1)$ for all $x_1,\, x_2 \in X$;  
    \item $c_2$ fulfilling the triangle inequality: $$c_2(x_1,x_2)\leq c_2(x_1,x_3)+ c_2(x_3,x_2)\quad \forall x_1,\, x_2, \, x_3 \in X.$$  
\end{enumerate}
Then, for every $0 \le n \tau = t$ and every $p \in \N$, we have
\[ c(\bar{x}_{t}^{\tau/2^p},\bar{x}_{t}^\tau)+c(\bar{x}_{t}^\tau,\bar{x}_{t}^{\tau/2^p}) \le \tau(2g(x_0)-g(\bar{x}_{t}^\tau)-g(\bar{x}_{t}^{\tau/2^p})).\]
    
\end{proposition}

\begin{proof}  From the nondegeneracy of $c$  coming from \eqref{eq:addass}  we get that $x_i^\tau=y_i^\tau$ for all $\tau \in (0,\bar{\tau})$ and all $i \in \N$. 
The assertion of the proposition follows from \eqref{eq:Cauchy_ineq_full}, as soon as we check that the quantity $A_c$  defined by 
\begin{multline}
    A_c:= \sum_{i=0}^{n-1}[c(\bar{x}_{i\tau}^{\tau/2^p},\bar{y}_{(i+1)\tau}^\tau)- c(\bar{x}_{(i+1)\tau}^{\tau/2^p},\bar{y}_{i\tau}^\tau) + c(\bar{x}_{i\tau}^\tau,\bar{y}_{(i+1)\tau}^{\tau/2^p})-c(\bar{x}_{(i+1)\tau}^\tau,\bar{y}_{i\tau}^{\tau/2^p})]\\
   \quad +2\sum_{i=0}^{n-1}c(\bar{x}_{(i+1)\tau}^\tau, \bar{y}_{i\tau}^\tau)+ 2\sum_{j=0}^{n2^p-1}c(\bar{x}_{(j+1)\tau/2^p}^{\tau/2^p}, \bar{y}_{j\tau/2^p}^{\tau/2^p})\label{eq:tph0}
\end{multline}
 is non-negative. 
Note that $A_c=A_{c_1}+A_{c_2}$. From the symmetry of $c_1\geq 0$ we obtain 
\begin{equation}
      A_{c_1}=2\sum_{i=0}^{n-1}c_1(\bar{x}_{(i+1)\tau}^\tau, \bar{y}_{i\tau}^\tau)+ 2\sum_{j=0}^{n2^p-1}c_1(\bar{x}_{(j+1)\tau/2^p}^{\tau/2^p}, \bar{y}_{j\tau/2^p}^{\tau/2^p}) \ge 0.\label{eq:tph1}
\end{equation}

Using the triangle inequality and the fact that $x_i^\tau = y_i^\tau $ and $x_i^{\tau/2^p} = y_i^{\tau/2^p} $, for  every $i=0, \dots, n-1$ one has that 
\begin{align*}
   c_2(\bar{x}_{(i+1)\tau}^\tau, \bar y_{i\tau}^\tau)+ c_2(\bar x_{i\tau}^\tau,\bar y_{(i+1)\tau}^{\tau/2^p})+c_2(\bar x_{(i+1)\tau}^{\tau/2^p}, \bar y_{i\tau}^{\tau/2^p}) &\ge c_2(\bar x_{(i+1)\tau}^\tau,\bar y_{i\tau}^{\tau/2^p}),\\
   c_2(\bar x_{(i+1)\tau}^{\tau/2^p}, \bar y_{i\tau}^{\tau/2^p})+c_2(\bar x_{i\tau}^{\tau/2^p},\bar y_{(i+1)\tau}^\tau) + c_2(\bar x_{(i+1)\tau}^\tau, \bar y_{i\tau}^\tau) &\ge c_2(\bar x_{(i+1)\tau}^{\tau/2^p},\bar y_{i\tau}^\tau).
\end{align*}
At the same time, again the triangle inequality gives 
\begin{equation*}
    2  \sum_{j=0}^{n2^p-1}c_2(\bar x_{(j+1)\tau/2^p}^{\tau/2^p}, \bar y_{j\tau/2^p}^{\tau/2^p}) \ge   2  \sum_{i=0}^{n-1}c_2(\bar x_{(i+1)\tau}^{\tau/2^p}, \bar y_{i\tau}^{\tau/2^p}).
\end{equation*}
By combining these inequalities one gets that 
\begin{multline}
   A_{c_2}\geq 
   \sum_{i=0}^{n-1}[c_2(\bar{x}_{i\tau}^{\tau/2^p},\bar{y}_{(i+1)\tau}^\tau)- c_2(\bar{x}_{(i+1)\tau}^{\tau/2^p},\bar{y}_{i\tau}^\tau) + c_2(\bar{x}_{i\tau}^\tau,\bar{y}_{(i+1)\tau}^{\tau/2^p})-c_2(\bar{x}_{(i+1)\tau}^\tau,\bar{y}_{i\tau}^{\tau/2^p})]\\
   +2\sum_{i=0}^{n-1}c_2(\bar{x}_{(i+1)\tau}^\tau, \bar{y}_{i\tau}^\tau)+ 2\sum_{i=0}^{n-1}c_2(\bar x_{(i+1)\tau}^{\tau/2^p}, \bar y_{i\tau}^{\tau/2^p}) \geq 0. 
   \label{eq:tph2}
\end{multline}
Inequalities \eqref{eq:tph1}--\eqref{eq:tph2} imply  that $A_c = A_{c_1} + A_{c_2} \geq 0$,  whence the thesis. 
\end{proof}

\end{document}
